% from ega3-6-fr.tex
\oldpage[III]{159}
Ce bifoncteur est l'aboutissement de six foncteurs spectraux biréguliers
\[
  {}^{(t)}\mathscr E
  (f_1,f_2;\mathscr P_\bullet^{(1)},\mathscr P_\bullet^{(2)})
  \qquad(t=a,b,a',b',c,d)
\]
dont les termes $E_2$ sont donnés par
\[
\begin{aligned}
  {}^{(a)}\mathscr E_{pq}^2
  &=\bigoplus_{q_1+q_2=q}
    \mathscr{T}\!or_p^S
    \bigl(\mathscr H^{-q_1}(f_1,\mathscr P_\bullet^{(1)}),
          \mathscr H^{-q_2}(f_2,\mathscr P_\bullet^{(2)})\bigr),\\
  {}^{(b)}\mathscr E_{pq}^2
  &=\mathscr H^{-p}\bigl(f_1\times_S f_2,
    \mathfrak{Tor}_q^S
    (\mathscr P_\bullet^{(1)},\mathscr P_\bullet^{(2)})\bigr),\\
  {}^{(a')}\mathscr E_{pq}^2
  &=\bigoplus_{q_1+q_2=q}
    \mathfrak{Tor}_p^S
    \bigl(\mathscr H^{-q_1}(f_1,\mathscr P_\bullet^{(1)}),
          \mathscr H^{-q_2}(f_2,\mathscr P_\bullet^{(2)})\bigr),\\
  {}^{(b')}\mathscr E_{pq}^2
  &=\mathscr H^{-p}\bigl(f_1\times_S f_2,
    \mathscr{T}\!or_q^S
    (\mathscr P_\bullet^{(1)},\mathscr P_\bullet^{(2)})\bigr),\\
  {}^{(c)}\mathscr E_{pq}^2
  &=\bigoplus_{q_1+q_2=q}
    \mathfrak{Tor}_p^S
    \bigl(f_1,f_2;
      \mathscr H_{q_1}(\mathscr P_\bullet^{(1)}),
      \mathscr H_{q_2}(\mathscr P_\bullet^{(2)})\bigr),\\
  {}^{(d)}\mathscr E_{pq}^2
  &=\mathscr H_p\bigl(\mathfrak{Tor}_q^S
    (f_1,f_2;\mathscr P_\bullet^{(1)},\mathscr P_\bullet^{(2)})\bigr).
\end{aligned}
\]
On dit que les suites spectrales $(a)$ et $(b)$ sont les \emph{suites
spectrales de Künneth}.

On notera que les suites spectrales $(a)$ et $(a')$ (resp. $(b)$ et $(b')$)
sont identiques lorsque $\mathscr P_\bullet^{(1)}$ et
$\mathscr P_\bullet^{(2)}$ se réduisent à leurs termes de degré 0; dans ce cas,
les suites $(c)$ et $(d)$ sont dégénérées et sont donc sans intérêt.
\end{theorem}

\begin{remark}[6.7.4]
Les hypertor globaux que nous avons définis ci-dessus comprennent comme cas
particuliers, à la fois les \emph{Modules d'hypercohomologie} définis dans
(6.2.1) et les \emph{hypertor locaux} définis dans (6.5.3). Montrons que l'on
a, pour tout morphisme $f:X\to Y$ \emph{quasi-compact et séparé} et tout
complexe $\mathscr P_\bullet$ de $\mathscr O_X$-Modules quasi-cohérents,
limité inférieurement, un isomorphisme canonique de $\partial$-foncteurs en
$\mathscr P_\bullet$
\[
  \mathfrak{Tor}_n^Y(f,1_Y;\mathscr P_\bullet,\mathscr O_Y)
  \xrightarrow{\sim}
  \mathscr H^{-n}(f,\mathscr P_\bullet)
  \qquad\text{(pour tout $n\in\mathbf Z$)}.
  \tag{6.7.4.1}\label{III.6.7.4.1.fr}
\]
En effet, les méthodes de recollement de (6.7.2) ramènent aussitôt au cas où
$Y$ est affine; on peut alors, en vertu de (6.2.2), calculer les deux membres
de (6.7.4.1) à l'aide d'un même recouvrement fini $\mathfrak U$ de $X$ par des
ouverts affines, et (pour le premier membre) du recouvrement de $Y$ formé de
$Y$ lui-même; avec les notations de (6.6.2), le bicomplexe
$L_{\bullet\bullet}^{(2)}$ est alors réduit à son terme de degrés $(0,0)$,
égal à $A$, et la conclusion résulte de ($0_{\mathrm I}$, 11.7.5). Pour une
généralisation de ce résultat, voir (6.7.7); mais on notera que lorsque dans
le premier membre de (6.7.4.1), on remplace $\mathscr O_Y$ par un
$\mathscr O_Y$-Module quasi-cohérent \emph{quelconque} $\mathscr F$, on n'a
plus en général un isomorphisme avec
$\mathscr H^{-n}(f,\mathscr P_\bullet\otimes_{\mathscr O_Y}\mathscr F)$, bien
que, dans le calcul précédent, le bicomplexe
$L_{\bullet\bullet}^{(1)}\otimes_A L_{\bullet\bullet}^{(2)}$ s'identifie
encore au bicomplexe
$C^\bullet(\mathfrak U,\mathscr P_\bullet\otimes_Y\mathscr F)$.

D'autre part, on a un isomorphisme canonique de bi-$\partial$-foncteurs
\[
  \mathfrak{Tor}_\bullet^S
  (1_{X^{(1)}},1_{X^{(2)}};
   \mathscr P_\bullet^{(1)},\mathscr P_\bullet^{(2)})
  \xrightarrow{\sim}
  \mathscr{T}\!or_\bullet^S
  (\mathscr P_\bullet^{(1)},\mathscr P_\bullet^{(2)}).
  \tag{6.7.4.2}\label{III.6.7.4.2.fr}
\]
En effet, on se ramène encore, par (6.7.1) et (6.7.2), au cas où $S$ et les
$X^{(i)}$ sont affines; dans le calcul du premier membre de (6.7.4.2), on peut
alors prendre pour
